\section{Empirical Evidence}
I document using matched employer-employee data that, wherever one observes higher layoffs, one also observes less responsive (to firm-level productivity shocks) wages, both across and within firms. Additionally, I document the connection between layoffs and labor production, average labor productivity, and wage dispersion. I will use these facts as testable implications of my model.

\subsection*{Data}
I use administrative data from France between 2009 and 2019. I combine a panel of workers from social security data containing 1/12th of the French labor force with annual data on firm balance sheet.  I focus on prime age workers (25-55 years old) and private jobs with wages $5\%$ above the national minimum wage at the time. Appendix~\ref{sample} provides more details on the sample selection. I end up with 265000 unique firms and 880000  unique workers per year.

I measure labor productivity using value added for worker, as provided in the balance sheet data. I model labor productivity $y_{fst}$ at firm $f$ in sector $s$ at time $y$ as 
\[ \log y_{fst} = \log a_t + \log b_{st}+ \log x_{fst}\]
where $a_t$ is the aggregate component, $b_{st}$ is a sectoral component and $x_{fst}$ is a firm-level component. I residualize $\log y_{fst}$ on time dummies to extract the common time component. I then measure the sectoral component $\log b_{st}$ as the average productivity across firms with a sector, and compute the firm component $\log x_{fst}$ as the residual. In the rest of the section, I will focus on the firm's responses to the firm-specific component $x_{fst}$ in order to abstract from any larger general equilibrium effects.

I measure wages as annual, CPI-adjusted labor earnings divided by the number of days worked. Labor earnings are net of payroll taxes but before income
taxes and they include all types of compensations, including bonuses and payment in kinds, but excludes stock options. I residualize the log of wages on occupation, firm, and region dummies as well as on a quadratic polynomial of worker experience. I remove the $5\%$ outliers from both the bottom and top of each year's distributions. \\
I measure layoffs as breaks in employment spells longer than 4 weeks. The intuition is that job-to-job transitions are usually do not entail long employment breaks, and, given low job-finding rate in France, it is unlikely for recently laid off workers to find employment within a month. There is still, however, risk of both false positives and negatives, as well as risk of misrepesenting voluntary worker movements into medium-term non-employment as layoffs. In Appendix~\ref{LFS} I use the French Labor Force Survey to measure layoffs as quarterly movements from employment into unemployment, as reported by the workers themselves.
\subsection*{Wages and Layoffs Across Firms}
I measure the response of wages as the percent change in wages after a $100\%$ increase in productivity, similarly to \textcite{souchier2022}. These responses are estimated using standard estimators from \textcite{guiso2005}. \\
I group firms based on their average layoff rates and measure the differences in response of log residualized wage growth to firm productivity shocks across these groups. I distribute firms into quantiles $d$ and compute the following response of wages:
\[\theta^{w,y}_d = \frac{Cov(\Delta \log(w_{ift}),\sum_{\tau=-1}^1 \Delta y_{ft+\tau})}{Cov(\Delta \log(y_{ift}),\sum_{\tau=-1}^1 \Delta y_{ft+\tau})}\]
where $y$ is the firm-level productivity, $\log(w_{ift})$ is the log residualized wage, $\sum_{\tau=-1}^1 \Delta y_{ft+\tau}$ is the 3-year sum of firm productivity growth. Although my model differs from that of \textcite{guiso2005}, I find it useful to report these statistics, instead of simple covariances, as they have been extensively studied in literature, and used for similar kind of empirics in papers close to mine.
%\[ \Delta \log w_{ift} = \sum_{d \in D} \mathbf{1}_{f \in d}(\alpha^d + \beta^{d}\Delta \log(x_{ft})) + \epsilon_{ift}\]
The estimates of wage passthrough across differently firing firms are shown in Table~\ref{tab:acrossfirms}. The results show that firms that lay off the largest share of their workers also exhibit the least responsive wages, all the way down to negatively responsive, consistent with the existing findings in the literature (\textcite{enrlich2024}). 

\begin{table} 
  \centering
  \begin{tabular}{lcc}
    \hline
    & \multicolumn{1}{c}{Layoff rate} & \multicolumn{1}{c}{Wage Change} \\   
    \hline
    \vspace{-4pt} 
    Layoff Gr.1     & $0.2\%$  &  $0.037^{***}$             \\    \vspace{-4pt} 
                    &    &  $(0.0009)$              \\    \vspace{-4pt} 
    Layoff Gr.2     & $2.5\%$  &  $0.036^{**}$              \\    \vspace{-4pt} 
                    &    &  $(0.0018)$              \\    \vspace{-4pt} 
    Layoff Gr.3    &  $6.4\%$  &  $0.020^{*}$                 \\    \vspace{-4pt} 
                    &    &  $(0.0021)$              \\    \vspace{-4pt} 
    Layoff Gr.4     & $14.6\%$  &  $0.019^{}$                \\    \vspace{-4pt} 
                    &    &  $(0.0020)$              \\                       
    Layoff Gr.5     & $33.9\%$  &  $-0.022^{}$                \\    \vspace{-4pt} 
                    &    &  $(0.0020)$              \\      
    \hline
  \end{tabular}
  \caption{Wage Passthrough across firms. Data: DADS Panel + FARE, 2009-2019}
  \label{tab:acrossfirms}
\end{table}

\subsection*{Wages and Layoffs Across Tenure} %Do I have this across both firms and tenure at the same time tho? I'm pretty sure I fucking don't...
%So then the question would be: what if the across firm heterogeneity is in fact just due to tenure differences? For that, I would need singificant composition differences in tenure across firms though.
%Regardless, kinda important to do these together, even if I don't interact them per se.
I introduce worker's tenure $ten \in T$ at the firm, provided to me directly in the employer-employee data. I look at wage passthrough and layoff rate vary with the first 5 years of tenure.
\[ EU_{ift} = \sum_{ten \in Ten} \mathbf{1}_{ift \in ten}\alpha^{ten}  + \epsilon_{ift}\] 
\[\theta^{w,y}_{ten} = \frac{Cov(\Delta \log(w_{ift}),\sum_{\tau=-1}^1 \Delta y_{ft+\tau})}{Cov(\Delta \log(y_{ift}),\sum_{\tau=-1}^1 \Delta y_{ft+\tau})}\]
Alongside the general wage passthrough, I look at its asymmetry. The core result of my model concerns what happens to wages right around layoffs, which would occur during negative productivity shocks. I look at how wages respond differently across tenure to the fact that the negative shock occured, rather than its actual value, in order to look at whether there is any asymmetry in the wage response, and how the asymmetry varies with tenure. \\
%\[ \Delta \log w_{ift} = \sum_{ten \in Ten} \mathbf{1}_{ift \in ten}(\alpha^{ten} + \beta^{ten}\Delta \log(x_{ft})) + \epsilon_{ift} \] 
The Table~\ref{tab:tenure} shows the results. I find that junior workers experience negative wage passthrough, that is, rising wages during negative productivity shocks, while more senior workers experience more positive higher wage passthrough. However, senior workers also experience much lower layoff rate (consistent with previous research, like \textcite{buhai2014}). This heterogeneity cannot be explained by standard stories of wage rigidity (minimum wage, sectoral bargainin, morale costs) without additional tweaks. Similarly, severance payments in France rise with tenure only barely. 

%\subsection*{Productivity and Wage Dispersion Response to Layoffs}
% Besides broader cross-sectional statistics computed above, my model has precise implications about what happens exactly when the firms lay off their workers. I will test these predictions by looking at the responses of labor productivity and wage dispersion to firm layoffs.
%\subsection*{Labor productivity response to layoffs} %When a firm fires, total production goes down, but per person goes up.
%Qualitatively, would that be true with the basic DRS? Or would the productivity stay the same? lol I feel stupid that I'm unsure
% No, that wouldn't be true in ANY case!!! 
% BUT WAIT!!! I was looking at the response to productivity shock y, not to layoffs. And, when y going down means L going up, we wouldn't see layoffs!
% In fact, in such a case, we would see layoffs when y goes up! and that might in fact be consistent with total prod going down, but avg going up.
% Generally ridiculous tho
%I use labor share within the firm as well as the per worker share. My model would predict that, upon layoffs, the labor share (which, in my model, is total production) would fall, but due to both quality improvements and downsizing, the per worker productivity will go up, in a similar vein to \textcite{berger2011}. To test that, I regress both measures of labor share on layoffs:
%\[ l_share_{f,t} = EU_{ft} + \epsilon_{ft}\]
\begin{table}
  \centering
  \begin{tabular}{lccc}
    \hline
    & \multicolumn{1}{c}{Layoff rate} & \multicolumn{1}{c}{Wage Passthrough} &  \multicolumn{1}{c}{Wage Response to Neg Shock}\\
    \hline
    \vspace{-4pt} 
    Less than 1 year        &  $21\%^{***}$   &  $-1.18^{***}$    \\    \vspace{-4pt} 
                        &  $(0.0002)$     &  $(0.0003)$     \\    \vspace{-4pt} 
    1-2 years        &  $6\%^{***}$       &   $-0.03^{***}$     \\     \vspace{-4pt} 
                        &  $(0.0002)$     &  $(0.0004)$       \\    \vspace{-4pt} 
    2-3 years        &  $4\%^{***}$       &  $-0.04^{***}$      \\     \vspace{-4pt} 
                        &  $(0.0002)$     &  $(0.0006)$       \\    \vspace{-4pt} 
    3-4 years        & $3\%^{***}$     &  $0.04^{***}$   \\ \vspace{-4pt}
                        &  $(0.0002)$     &    $(0.0007)$ \\    \vspace{-4pt} 
    4-5 years        & $2\%^{***}$     & $0.06^{***}$  \\    \vspace{-4pt} 
                        &  $(0.0002)$     &    $(0.0009)$ \\  

    \hline
  \end{tabular}
  \caption{Layoffs and Wages passthrough across tenure. Data: DADS Panel + FARE, 2003-2019}
  \label{tab:tenure}
\end{table}
\subsection*{Senior/junior wage ratio response to layoffs} %When a firm fires, wage differential across tenure should go down
Following \ref{prop:qualityconv}, the model predicts that junior workers will be the first to get cut during a layoff round. Then, following \ref{proof:wage_cuts}, it predicts that the remaining workers will higher wage growth. The immeadiate implication of this is that, following layoffs, we should expect the wage dispersion across cohorts to contract. To check this in the data, I look at how the wage ratio between senior and junior workers responds to layoffs. For each firm-year, I compute the median tenure, and take the ratio of wages of workers above and below the median $\bar{w}_{sen,f,t}/\bar{w}_{jun,f,t}$. I then compute log difference of this ratio across time and regress it on layoffs:
\[ \Delta \bar{w}_{sen,f,t}/\bar{w}_{jun,f,t} = EU_{ft} + \epsilon_{ft} \]

The estimate is shown in Figure~\ref{tab:layoff_response}. I find that both implications of my model hold qualitatively in the data: labor share falls with layoffs, per worker labor share rises, and wage dispersion across tenure falls.

\begin{table}
  \centering
  \begin{tabular}{lcc}
    \hline
    & \multicolumn{1}{c}{Layoff} & \multicolumn{1}{c}{Wage Change} \\
    \hline
    \vspace{-4pt} 
    Size Gr.1        &  $0.07^{***}$   &      \\    \vspace{-4pt} 
                        &  $(0.0002)$     &       \\    \vspace{-4pt} 
    Firm Shock * Size Gr. 3        & $-0.009^{***}$     &   $0.003^{***}$ \\    \vspace{-4pt} 
                                    &  $(0.0001)$     &    $(0.0003)$ \\  
    \hline
  \end{tabular}
  \caption{Layoffs and Wages passthrough across firm size. Data: DADS Panel + FARE, 2003-2019. \textcolor{red}{TEMPORARY FIGURE. TO BE OVERHAULED}}
  \label{tab:layoff_response}
\end{table}
